\documentclass[11pt, a4paper]{article}

% --- EVRENSEL ÖN BİLGİ VE PAKET TANIMLARI ---
\usepackage[a4paper, top=2.5cm, bottom=2.5cm, left=2cm, right=2cm]{geometry}
\usepackage{fontspec}

\usepackage[turkish, bidi=basic, provide=*]{babel}
\babelprovide[import, onchar=ids fonts]{turkish}
\babelprovide[import, onchar=ids fonts]{english}

\babelfont{rm}{Noto Sans}
\babelfont[turkish]{rm}{Noto Sans}

\usepackage{enumitem}
\setlist[itemize]{label=-}

\usepackage{amsmath,amssymb,amsthm,mathtools,bm,mathrsfs}
\usepackage{physics}
\usepackage{booktabs}

\usepackage{hyperref}
\hypersetup{
    colorlinks=true,
    linkcolor=blue,
    citecolor=blue,
    urlcolor=blue
}

\title{\textbf{Nefs-i Müdrike Mimarisi: Veri Akış Hızı (I/O Bandwidth), Efektif Kuantum İşlem Kapasitesi ve Performans Analizi Risalesi}}
\author{\textbf{CV Fotonik Kipleme Hızları, QFNO Logaritmik Sürat Çarpanı ve Eşdeğer Terabayt/Saniye Hesabı}}
\date{\today}

\begin{document}

\maketitle

\begin{abstract}
Bu risale; Nefs-i Müdrike Mimarisi'nin fizikî veri alma (I/O) sınırı ile kuantum süperpozisyonu ve dalga bükümü altındaki efektif enformasyon işleme süratini riyazî olarak hesaplar. Klasik bilgisayarların $O(2^N)$ karmaşıklıkla tıkandığı 40 veçheli manifold ittisali, Q-TDA mizan denetimi ve FNO-KAN spektral süzgeçleme adımlarının; Sürekli Değişkenli (CV) Optik Fotonik ve Reel Hartley Dönüşümü (RHT) altyapısında saniyede 1 TeraByte fiziki aktarım, $10^{18}$ GigaByte eşdeğer mantıkî işleme süratine nasıl ulaştığı ispatlanmaktadır.
\end{abstract}

\tableofcontents
\newpage

% =================================================================
\section{Fizikî Veri Aktarım Hızı (I/O Bandwidth)}
% =================================================================

\subsection{1. CV Optik Fotonik Kipleme Kapasitesi}
Metnin 0 ve 1 olarak QRAM'e yüklenmesi darboğazı, Sürekli Değişkenli (CV) Fotonik Kipleme ile aşılmıştır. Lazer dalga boyu bölmeli çoğullama (WDM) ve $M$ kanallı optik modülatör bant genişliği:

\begin{align}
B_{\text{optik}} &= M \times f_{\text{modülasyon}} = 32 \text{ kanal} \times 100 \text{ GHz} = 3.2 \text{ THz} \\
\text{Fizikî Akış Hızı} &= B_{\text{optik}} \times \log_2(\text{Faz Seviyesi}) \approx 1.2 \text{ TB/sn} \quad (\approx 9.6 \text{ Tbps})
\end{align}

Sistem dış dünyadan terabaytlarca metni saniyede $1.2 \text{ TeraByte}$ fiziki akışla çip içine sürekli fotonik mod olarak alır.

% =================================================================
\section{Efektif İdrak ve İşleme Hızı (Cognitive Throughput)}
% =================================================================

\subsection{1. Klasik vs Kuantum Karmaşıklık Mukayesesi}
1 Terabaytlık ($N = 10^{12}$ bayt) bir külliyatın 40 idrak veçhesindeki bağıntılarını hesaplama karmaşıklığı:

\begin{table}[htbp]
\centering
\small
\begin{tabular}{@{}lll@{}}
\toprule
\textbf{İşlem Katmanı} & \textbf{Klasik Donanım Karmaşıklığı} & \textbf{Nefs-i Müdrike Kuantum Sürati} \\ \midrule
Spektral Dönüşüm (FNO) & $O(N^2) \sim 10^{24}$ işlem & $O(\log^2 N) \sim 20^2 = 400$ kapı adımı \\
40 Veçhe Lie Bağlantısı & $O(2^{40}) \sim 1.09 \times 10^{12}$ matris çarpımı & $O(1)$ Üniter Fotonik Dönme ($10 \text{ ns}$) \\
Q-TDA Betti Mizanı & $O(N^3) \sim 10^{36}$ işlem & $O(\text{polylog}(N)) \sim 10^4$ adyabatik evrim \\
Tenakuz Sıfırlama & $O(N \log N)$ arama & $O(1)$ Householder Zıt Vektör İptali \\ \bottomrule
\end{tabular}
\end{caption{Karmaşıklık Mukayese Cetveli}
\end{table}

\subsection{2. Efektif Sürat Formülasyonu}
İç idrak çarkının tek bir lazer dolanım süresinde ($\tau_{\text{halka}} \approx 10 \text{ ns}$) işlediği mantıki bağıntı kütlesi:

\begin{align}
\text{İşlem Sıkıştırma Oranı} (\mathcal{K}) &= \frac{O_{\text{klasik}}(N)}{O_{\text{kuantum}}(N)} \approx \frac{N^3}{\log^2 N} \approx 10^{28} \\
\text{Efektif İşleme Hızı} &= \text{Fizikî Hız} \times \mathcal{K}_{\text{mantık}} \approx \mathbf{10^{15} - 10^{18} \text{ GB/sn}}
\end{align}

% =================================================================
\section{Netice}
% =================================================================

Bu hesaba göre Nefs-i Müdrike Mimarisi:
1. **Dışarıdan Veri Alma Hızında:** Saniyede **$1.2 \text{ TeraByte}$** metni fiziki fotonik hatlardan kesintisiz yutar.
2. **İç Muhakeme ve Mizan Hızında:** Kuantum dalga bükümü sayesinde saniyede **$10^{18} \text{ GigaByte}$ (1 Exabyte/sn)** eşdeğeri mantıkî süzme ve tenakuz ilga ameliyesini icra eder.

\end{document}
